\subsubsection{4.2 可行性分析}

{\kaishu
	本项目\textbf{面向国家重大应用需求，目标明确，思路清晰}，并已具备\textbf{坚实的技术基础和良好的交流合作}，因而\textbf{整体项目实施切实可行}，具体如下：
	
	\textbf{（1）需求迫切}
	
	当前，海洋安全、低空安全和无人系统监管需求快速增长，海上无人艇、低空无人机、漂浮小目标等新型无人目标呈现小型化、低慢化、机动化和隐蔽化特征，对电磁感知侦测系统的实时发现、稳定跟踪和智能识别能力提出了更高要求。海洋环境中，海面起伏、多径传播、风浪扰动和复杂电磁背景共同作用，使海杂波呈现强非平稳、强时变、宽动态和重尾分布等特征，微弱目标信号极易被强海杂波淹没，导致传统基于固定门限、静态模型和离线处理的频谱分析方法难以适应复杂开放海况下的无人目标侦测需求。特别是在远海、岛礁、港口、航道巡检和无人平台协同感知等应用场景中，前端设备通常面临算力、存储和通信资源受限的问题，难以对连续产生的大规模电磁频谱数据进行全量保存和集中处理，进一步制约了弱小目标信号的及时提取与有效利用。
	
	因此，迫切需要开展海杂波场景下电磁频谱数据智能处理机理与方法研究，揭示复杂海况下频谱信号传播与杂波干扰规律，突破海杂波干扰下微弱信号提取、有限算力端侧频谱数据高效存储、边缘频谱数据即时智能处理等关键技术，形成面向复杂海洋环境的“感知—存储—处理—验证”一体化方法体系。本项目的开展不仅有助于提升海杂波背景下无人目标电磁感知侦测的可靠性和实时性，也将为海洋安全监测、低空与海空协同防控、无人平台智能感知和复杂环境下电磁频谱资源利用提供重要技术支撑，具有突出的现实需求和战略意义。
	
	\textbf{（2）目标明确}
	
	xxx
	
	\textbf{（3）思路清晰}
	
	xxx
	
	
	\textbf{（4）技术扎实}
	
	项目组在前期科研工作中围绕电磁频谱感知、海杂波建模与目标检测、人工智能理论与应用、大规模数据管理、端侧存储优化、分布式可扩展存储、可信数据管理和智能信息处理等方向积累了较为扎实的研究基础，形成了一系列与本项目任务密切相关的理论方法、算法模型和系统研发经验。项目组成员长期从事雷达信号处理、复杂环境目标探测、人工智能、知识图谱、深度学习、跨模态检索、大数据存储、数据库软硬协同优化、新型存储机制、高并发可扩展存储、分布式数据管理、可信存储与安全传输等研究，研究方向覆盖本项目所需的海杂波场景电磁频谱信号传播机理分析、微弱信号提取、有限算力端侧频谱数据存储、频谱知识库构建与动态管理、端边协同频谱数据可信传输以及边缘频谱数据即时智能处理等关键环节。
	
	在前期工作基础方面，项目组已在人工智能、数据挖掘、知识图谱、智能检索、大数据存储、数据库系统、软硬件协同优化、分布式可扩展存储、可信数据管理和复杂信号智能处理等领域重要期刊和会议（IEEE Transactions on Affective Computing、ACM Transactions on Software Engineering and Methodology、IEEE TCAD、ACM TACO、PVLDB、AAAI、NeurIPS、IJCAI、DATE、ICPP、ICCD、ICDCS）发表多篇高水平论文，并承担国家级、省部级和横向合作项目，形成了较好的系统研发与工程验证基础。此外，项目组已围绕本项目研究方案开展多次研讨，\textcolor{red}{对海杂波场景频谱数据智能处理、端侧高效存储、频谱知识库动态管理和可信传输等关键问题形成了明确方向和清晰思路。}因此，\textbf{本项目研究具有较高的技术可行性。}
	
	\textbf{（5）队伍合理}
	
	本项目团队成员由\textbf{武汉理工大学}、\textbf{东北大学}和\textbf{联通数字科技有限公司}的相关科研人员组成，其中教授1人（郑渤龙）、副教授2人（王亮亮、邓诗卓）、青年教师2人（董明、王海鑫）、高级工程师2人（魏春城、蒋维）、工程师1人（张少辉）。成员主体都是科研一线的中青年学者，近年来在国际重要期刊和会议上发表100余篇高水平论文，并承担多项国家自然科学基金和省部级科研项目，已在学术界形成一定影响力。团队优势互补：武汉理工大学团队长期致力于人工智能、具身智能、分布式数据管理等研究，东北大学长期致力于具身智能、机器人视觉等研究，联通数字科技有限公司长期致力于“云物融合架构下的具身智能”技术创新与应用，具有较好的工程实践能力。团队成员分工明确，能够保证充足的时间与精力投入。综上，\textbf{本项目的团队结构合理，可以为本项目的顺利实施提供优秀的人力保障}。
	
	\textbf{（6）合作良好}
	
	\textcolor{red}{本项目申请人已与多位国（境）外同行专家学者保持密切的学术交流与合作，主要包括丹麦奥尔堡大学Christian S. Jensen教授（欧洲科学院院士、丹麦皇家科学院院士、ACM/IEEE Fellow）、香港科技大学Xiaofang Zhou教授（IEEE Fellow）、香港中文大学Yufei Tao教授（ ACM Fellow、SIGMOD 2013、2015和PODS 2018等最佳论文奖获得者）、新加坡国立大学Xiaokui Xiao教授（IEEE Fellow、VLDB 2021最佳研究论文奖和SIGMOD 2022研究亮点奖等获得者）等。项目执行期间，拟进一步加强与这些专家学者的交流合作，\textbf{这也将为本项目研究的顺利展开提供强有力的技术支持和良好的学习研讨氛围。}}
}